\chapter{Retinal Detachment Symptoms}

\section*{Definition}
Retinal detachment (RD) is the separation of the neurosensory retina from the underlying retinal pigment epithelium (RPE), resulting in photoreceptor ischemia and rapid vision loss if not repaired within hours to days.

\section*{What Happens in Your Body}
A full-thickness retinal break (riegmatogenous RD) allows liquefied vitreous to pass through the defect and accumulate in the subretinal space, peeling the neurosensory retina off the RPE. Tractive forces from vitreoretinal adhesions or exudative fluid accumulation may also cause detachment.

\section*{Causes}
\begin{itemize}
  \item \textbf{Rhegmatogenous (most common):} Retinal tear associated with PVD, high myopia, lattice degeneration, or trauma
  \item \textbf{Tractional:} Fibrovascular membranes in proliferative diabetic retinopathy or retinopathy of prematurity
  \item \textbf{Exudative:} Inflammatory (uveitis, Vogt-Koyanagi-Harada), neoplastic (choroidal melanoma), or vascular (Coats disease)
  \item \textbf{Risk factors:} Previous cataract surgery, prior RD in fellow eye, blunt ocular trauma, high myopia, and family history
\end{itemize}

\section*{When to See a Doctor}
See your doctor if:
\begin{itemize}
  \item Sudden onset of bright flashing lights (photopsia) in the temporal field
  \item Rapid increase in number of floaters described as a "shower of soot" or cobwebs
  \item Curtain or shadow progressing across the visual field
  \item Sudden, painless, unilateral vision loss with a relative afferent pupillary defect
  \item Any new photopsia or floater in a high-risk patient (myope, post-cataract surgery, prior RD)
\end{itemize}

\section*{Related Symptoms}
\begin{itemize}
  \item Photopsia (flashing lights)
  \item Floaters (sudden increase, often described as black dots or cobwebs)
  \item Visual field defect (curtain-like, progressing centrally)
  \item Central vision loss once macula detaches
  \item Metamorphopsia (distortion)
\end{itemize}
\textbf{Important:} Central vision is salvageable if the macula remains attached at time of repair; prognosis declines sharply if the macula is detached for more than 72 hours.

\section*{Self-Care / Home Management}
\begin{itemize}
  \item Strict head positioning to keep the dependent retinal break in the most gravity-dependent location (e.g., face-down for superior breaks)
  \item Avoidance of strenuous activity, heavy lifting, and rapid head movements
  \item Urgent ophthalmology referral -- do not delay for any home management
  \item Shield the unaffected eye to protect it from injury (especially in monocular people)
  \item No administration of any eyedrops until examined
\end{itemize}

\section*{How Your Doctor Will Evaluate You}
\begin{itemize}
  \item Dilated fundus examination with binocular indirect ophthalmoscopy and scleral depression
  \item Slit-lamp biomicroscopy with a 90D or 78D lens for macular assessment
  \item B-scan ultrasonography if media opacity (cataract, vitreous hemorrhage) prevents retinal visualization
  \item Goldmann three-mirror examination of the peripheral retina
  \item Optical coherence tomography (OCT) to confirm shallow macular involvement
\end{itemize}

\section*{Treatment Options}
\begin{itemize}
  \item \textbf{Pneumatic retinopexy:} Intravitreal gas injection with cryotherapy or laser to the break, combined with strict head positioning (selected superior breaks)
  \item \textbf{Scleral buckle:} Silicone band or sponge encircling the globe to indent the wall and close the break
  \item \textbf{Pars plana vitrectomy:} Removal of vitreous, laser retinopexy, and gas or silicone oil tamponade
  \item \textbf{Combined approach:} Vitrectomy with scleral buckle for complex or inferior detachments
  \item Postoperative positioning (face-down or other orientation as directed by the surgeon)
  \textit{} Avoid air travel and high altitude until gas bubble fully resorbs
\end{itemize}

\section*{References}
\begin{thebibliography}{99}
\bibitem{retinal_detachment_symptoms:ref1} D'Amico DJ. "Clinical practice: Primary retinal detachment." \textit{New England Journal of Medicine}, 2008;359(22):2346--2354.
\bibitem{retinal_detachment_symptoms:ref2} Schwartz SG, Flynn HW Jr, Mieler WF. "Update on retinal detachment surgery." \textit{Current Opinion in Ophthalmology}, 2016;27(3):175--181.
\bibitem{retinal_detachment_symptoms:ref3} Mitry D, Charteris DG, Fleck BW, et al. "The epidemiology of rhegmatogenous retinal detachment: Geographical variation and clinical associations." \textit{British Journal of Ophthalmology}, 2010;94(6):678--684.
\bibitem{retinal_detachment_symptoms:ref4} Boberg-Ans G, Henning V, Villumsen J, et al. "Long-term outcomes after primary rhegmatogenous retinal detachment repair." \textit{Acta Ophthalmologica}, 2020;98(2):e214--e221.
\bibitem{retinal_detachment_symptoms:ref5} Steel DHW, Lotery AJ, UK Retinal Detachment Outcomes Study Group. "Factors influencing outcomes in rhegmatogenous retinal detachment." \textit{Eye}, 2021;35(10):2701--2710.
\bibitem{retinal_detachment_symptoms:ref6} Yanoff M, Duker JS. \textit{Ophthalmology}, 5th ed. Elsevier, 2019.
\end{thebibliography}
